% Fate's Edge SRD --- Advancement & XP (Section 15)

\section{Advancement \& XP}

Advancement in Fate's Edge reflects meaningful growth in capability and standing. XP is awarded for table-facing accomplishments, hard choices, and the dramatic friction that defines your story. Boons can be converted sparingly to accelerate growth.

\subsection{Awarding XP}

Choose a session pacing dial and stick to it for a campaign arc.

\begin{itemize}
  \item \textbf{Gritty:} 4--6 XP per session (slow burn).
  \item \textbf{Standard:} 6--10 XP per session (default pace).
  \item \textbf{Heroic:} 10--14 XP per session (fast growth).
\end{itemize}

\subsubsection{Session Awards (Guidelines)}

\begin{tabular}{@{}lp{5cm}@{}}
\toprule
\textbf{Award} & \textbf{XP} \\
\midrule
Table attendance & +2 \\
Major objective reached & +2--4 \\
Discovery or lore unlocked & +1--2 \\
Hard choice embraced & +1--2 \\
Complication spotlight & +1--3 \\
Bond or flag driven play & +1--2 \\
GM curveball award & +0--3 \\
\bottomrule
\end{tabular}

\subsubsection{Milestones}
\begin{itemize}
  \item At the conclusion of a major story arc, award +8--12 XP to all players.
  \item Grant +2 XP to one player for a signature moment of the arc.
\end{itemize}

\subsubsection{Boon Conversion}
Once per session, during downtime, a character may convert 2 Boons $\rightarrow$ 1 XP (max 2 XP via conversion per session). All normal Boon limits apply (hold 5; trim to 2 at scene end).

\begin{quote}
\emph{Design Note:} Boon conversion is a trickle, not a primary advancement path. The real value of Boons is in the moment—re-rolling a failure, activating an Asset, or improving Position. Hoarding Boons for conversion starves the table of drama. Spend them freely; the XP will come from your adventures.
\end{quote}

\subsection{Spending XP}

\begin{tabular}{@{}lp{3.8cm}p{5cm}@{}}
\toprule
\textbf{Improvement} & \textbf{Cost} & \textbf{Downtime} \\
\midrule
Attribute increase & new rating $\times$ 3 XP & new rating (days) \\
Skill increase & new level $\times$ 2 XP & new level (days) \\
On-screen Follower & Cap$^2$ XP & 1--3 days (recruit and brief) \\
Minor Asset & 4 XP & 1 day \\
Standard Asset & 8 XP & 1 week \\
Major Asset & 12 XP & 1 month \\
Talent & as listed (2--12+ XP) & as listed \\
\bottomrule
\end{tabular}

\subsection{Upkeep}

Upkeep is paid once per Downtime period per asset or follower.

\begin{itemize}
  \item \textbf{Efficient (Higher XP, Less Time):} Upkeep XP = $\max(1, \lceil \frac{\text{Acquisition XP}}{3} \rceil)$. Minimal effort.
  \item \textbf{Intensive (Lower XP, More Time):} Cost = 1 XP. Requires a dedicated Downtime action with significant personal involvement.
\end{itemize}

\textbf{Failure to pay:} The resource degrades.
\begin{itemize}
  \item \textbf{Follower:} Loyalty worsens one step (Faithful $\rightarrow$ Strained $\rightarrow$ Broken).
  \item \textbf{Asset:} Status worsens one step (Flourishing $\rightarrow$ Stable $\rightarrow$ Strained $\rightarrow$ Collapsed).
\end{itemize}

\subsection{Rush Rule}

When you buy or upgrade an ability during downtime, you may skip the normal training time by accepting a risk.

\subsubsection{Rush Rule Procedure}

\begin{enumerate}
  \item \textbf{Declare:} The player states they are rushing their training.
  \item \textbf{Set Timer:} The GM creates a \textbf{Haste Timer [4]}.
  \item \textbf{Push Forward:} Each downtime action during the rush, the player rolls \textbf{Spirit + Endurance} (DV 3) to represent pushing through fatigue and distraction.
    \begin{itemize}
      \item \textbf{Clean Success:} No tick. Progress continues smoothly.
      \item \textbf{Partial:} The timer ticks +1. The training is harder than expected.
      \item \textbf{Miss:} The timer ticks +2. Something goes wrong—a pulled muscle, a broken tool, a moment of doubt.
    \end{itemize}
  \item \textbf{Complete:} If the timer reaches 4 before training completes, the ability or asset is flawed.
    \begin{itemize}
      \item \textbf{Attribute/Skill:} The new rating is gained, but the character suffers a temporary -1 die penalty until they spend a full downtime session retraining (removing the flaw).
      \item \textbf{Asset:} The asset is acquired but has a hidden complication (GM chooses: a rival knows about it, it requires unusual upkeep, or it's in a bad location).
      \item \textbf{Talent:} The talent works, but the GM gains 2 SB to spend on a complication tied to its rushed nature (e.g., the technique is flashy but draws attention, or it leaves an opening).
    \end{itemize}
  \item \textbf{Clean Success:} If the timer never ticks (0 ticks after all required downtime actions), the ability is learned without flaw.
\end{enumerate}

\subsection{Tiers of Reputation}

Reputation tiers reflect how the world responds to you and what your character can accomplish.

\begin{tabular}{@{}llll@{}}
\toprule
\textbf{Tier} & \textbf{XP Range} & \textbf{Standing} & \textbf{Mechanical Effect} \\
\midrule
I --- Rookie & 0--40 & Local & Max dice pool 8d10; max 2 Followers; Prestige talents locked \\
II --- Seasoned & 41--90 & Regional & Max dice pool 9d10; max 3 Followers; Prestige talents (7--10 XP) unlocked \\
III --- Veteran & 91--150 & National & Max dice pool 10d10; max 4 Followers; Epic talents (11+ XP) unlocked \\
IV --- Paragon & 151--220 & Continental & +1 max Boon hold (6); Factions respond to your requests \\
V --- Mythic & 221+ & Legendary & +1 max Boon hold (7); World-shaping narrative permissions \\
\bottomrule
\end{tabular}

\subsection{High-Tier Play and Boons}

As characters advance, the role of Boons shifts.
\begin{itemize}
  \item At early tiers, Boons primarily buy reliability --- negating Fatigue, stabilizing risky actions, or nudging Position.
  \item At higher tiers, Boons increasingly buy authority. A single Boon may force a truth, seal a pact, redirect consequences, or resolve a phase of a conflict outright.
\end{itemize}

High-tier characters typically have fewer Boons to rely on, but each spend carries greater narrative weight. This is intentional. Power at high tier comes from preparation, Bonds, oaths, and hard choices --- not from inflating Difficulty Values. The world does not become harder to roll against simply because characters are stronger; instead, success and failure both reshape the story in larger ways.

\subsubsection{High-Tier Example: The Archmage's Choice}

A Tier IV archmage attempts to counterspell a world-ending ritual. She has 0 Boons left. She rolls 8 dice (Wits 5 + Arcana 3) vs DV 4. Result: 3 successes — a Partial.

\textbf{Without Boons:} She cannot re-roll. The GM rules: "The ritual slows, but a backlash rips through the ley lines. A nearby village is consumed by wild magic. You save the world, but at a cost."

\textbf{With a Boon:} She could re-roll a failure and possibly achieve a Clean Success, avoiding the backlash.

The scarcity of Boons made the Partial outcome feel meaningful, not like a failure. At high tiers, every Boon is precious—spend them wisely.

\subsection{Advancement Notes}

\begin{itemize}
  \item Attribute cost scaling encourages diversification over single-stat spikes.
  \item Skill mastery yields tangible, fiction-first benefits.
  \item Prestige abilities (6+ XP) should be gated by narrative milestones or patron bargains.
\end{itemize}